\section{A sharp source bound for the second positive forward difference}
Let $F=\sum f_nz^n$ be an actual nonrectangular increasing arm.
Set $D_0=D_1=0$, $D_n=(f_n-f_{n-1})^+$ for $n\ge2$, and
\[
 W_n=(D_n-D_{n-1})^+,\qquad S_K=z+\cdots+z^K.
\]
Write $\kappa\ge3$ for the first missing weight and $K=\kappa-1$.
All sequences have zero extension. The following statement also
holds for formal append words whose first missing weights never decrease.

\begin{theorem}[Second-difference source]
Coefficientwise on the whole polynomial,
\[
                  (1-S_K)W\preceq z^2-z^\kappa.       \tag{1}
\]
In particular the negative unit at the first missing weight is
retained; it is not replaced by a bound on cumulative sums.
For every fixed $K\ge2$ equality holds on every prescribed finite
range for a sufficiently long actual rectangular host $K^L$.
\end{theorem}

\begin{proof}
We use the proved ordinary positive-difference shape theorem:
$D$ is unimodal, starts at rank two with value one, and has
initial interval support. Let $m$ be its last mode. The ordinary
positive-difference splice for a legal append $t$ says that the
new $D'$ agrees initially with
\[
                   P=z^2[t-1]+z[t]D,
\]
through one rank after the old ordinary last mode, and is
nonincreasing afterwards. Put
\[
                   R=z^2-z^{t+1}+S_tW.                \tag{2}
\]
This is nonnegative: its only possible subtraction is at $t+1$,
where the $W_2=1$ term occurs in $S_tW$.

First prove the full source bound for $R$. There are exactly two
legal first-missing phases. In the stationary phase $t=K$,
the old and new $K$ agree, and
\[
 (1-S_K)R\preceq(1-S_K)(z^2-z^{K+1})
                         +S_K(z^2-z^{K+1})
                    =z^2-z^{K+1}.                    \tag{3}
\]
In the growth phase write $K=\kappa_{\rm old}$, so the new
first missing weight is $K+1$, $t\ge K$, and $K\ge3$.
The old source bound, now with the extra summand $z^K$, is
\[
              (1-S_K)W\preceq z^2-z^K-z^KW.
\]
After substitution in (2), the difference between the new target
$z^2-z^{K+1}$ and this upper bound for $(1-S_K)R$ is
\[
 \begin{split}
 &z^{K+1}[t](W-z^2)\\
 &\quad+z^{K+2}\bigl([t-1]-z^{t-K}\bigr)
                 +z^{t+1}+z^{t+K+2}.                 \tag{4}
 \end{split}
\]
Every term is nonnegative: $W\succeq z^2$ and
$0\le t-K\le t-2$. This proves the raw source inequality.

We now check that replacing $R$ by the actual $W'$ preserves it.
For every $n\le m+1$, all sampled old differences of $D$ are
nonnegative, and $W'_n=R_n$. At $n=m+2$ the new value is at
most $R_n$, while every earlier value still agrees. The source
inequality therefore holds through this first crossing edge.

After that edge, before the final ordinary splice, the underlying
values
\[
        v_n=(D_{n-1}-D_{n-t-1})^+
\]
are nonincreasing. Indeed the entering index is after the old
$D$ mode while the leaving index is before it; once both are
after the mode the difference is nonpositive. The only exception
in $W'$ is a reduction by one at $n=t+1$. Thus every later source
comparison is paid by its immediately preceding term, except
possibly $n=t+2$. At this exceptional recovery,
$W'_{t+2}\le W'_{t+1}+1$. If $W'_t=0$, monotonicity of $v$
forces $W'_{t+2}=0$; otherwise $W'_t\ge1$ pays the extra unit.
Since $K\ge2$, these two preceding terms are both available.
The final ordinary splice can only remove the remaining positive
second differences, so it introduces no new recovery.

These terminal arguments prove the bound zero at ranks other
than two and $\kappa'$. At $\kappa'$ the required negative
unit follows directly from actual universality: below $\kappa'$
the ordinary ranks are full. If $W'_{\kappa'}>0$, its value is
the full value minus the first defect, at least one. If it is
zero, its preceding window still contains $W'_2=1$.
In either case the source is at most minus one.

For a rectangular old arm, $D'$ is an interval of ones and
$W'=z^2$, with new $\kappa'=3$, establishing the base.
This proves induction. For sharpness a sufficiently long $K^L$
agrees in any prescribed range with $(1-S_K)^{-1}$. Its positive
differences and second positive differences there satisfy
\[
 D=\frac{z^2[K-1]}{1-S_K},\qquad
 W=\frac{z^2-z^{K+1}}{1-S_K},
\]
giving the claimed equality.
\end{proof}

\begin{corollary}[A delayed local curvature payment]
Whenever $n\ge3$ and $D_n>D_{n-1}$,
\[
          2D_{n-1}-D_n\ge {\bf1}_{n=\kappa}+D_{n-\kappa}.
                                                               \tag{5}
\]
In particular, for every $n\ge\kappa+2$,
\[
       2^{\kappa-1}D_n
          \le(2^\kappa-1)D_{n-1}.                      \tag{6}
\]
The strict growth factor in (6) is a consequence, not a claim
of the sharp exponential constant for the whole family.
\end{corollary}

\begin{proof}
Let $V=1-(1-S_K)F$. Since $D$ is unimodal, the hypothesis
$D_n>D_{n-1}$ puts every earlier difference of $D$ on its
nonnegative prefix. On this prefix
$(1-z)^2F=1-z+W$.
The identity $(1-z)(1-S_K)=1-2z+z^\kappa$ and (1) imply
$\Delta^2V_n\ge0$ for $n\ge3$.
As $E=(1-z)V+z^\kappa F$, taking a difference gives (5).
For $n\ge\kappa+2$, backward iteration of the established
single-arm doubling bound gives
$2^{\kappa-1}D_{n-\kappa}\ge D_{n-1}$; combine with (5).
If $D_n\le D_{n-1}$, (6) is immediate. This proves all cases.
\end{proof}

This is an additional actual-arm restriction. It does not assert
log-concavity or unimodality of $W$, either of which would be
stronger and is not used. General property C, the full later mixed
budgets and the original unamplified conjecture remain OPEN.
